\documentclass[10pt]{article}
\usepackage[margin=0.8in]{geometry}
\usepackage{booktabs}
\usepackage{hyperref}
\usepackage{amsmath}
\title{Active Cost-Sensitive Monitoring of Black-Box LLM Behavior}
\author{Week 1 Project Prototype}
\date{}
\begin{document}
\maketitle
\begin{abstract}
We design a small monitoring agent for externally observable language-model behavior when the true model state is hidden. The agent maintains a belief over stable behavior, persistent degradation, transient anomalies, and distribution shift, then selects ACCEPT, INVESTIGATE, or REJECT using an explicit cost model. A deterministic 50-case simulation compares a static threshold baseline, a binary belief policy, and an active policy that can purchase additional evidence under a bounded investigation budget. The experiment is a reproducible engineering test, not a production deployment claim.
\end{abstract}

\section{Introduction}
Hosted LLM services can change behavior over time even when the service appears to expose the same model family, motivating continuous external evaluation \cite{chen2023llmdrift}. The Week 1 problem is therefore: the agent observes external LLM inputs, outputs, evaluation signals, and runtime evidence and must choose an action because the provider-side state is not directly visible.

\section{Related Work and User Discussions}
The LLMDrift work provides historical generations and evaluation evidence for studying behavior changes over time \cite{chen2023llmdrift}. Public discussions for this project also identified three engineering constraints: the monitor cannot infer internal model state, baseline scores must be computed consistently, and investigation must have a stopping budget.

\section{Agent Design}
Observable evidence includes task quality, semantic drift, formatting/safety changes, latency, and error rate. Hidden states are $S=\{stable,degraded,transient,distribution\_shift\}$. The belief vector $b(s)$ is updated by a naive-Bayes-style likelihood model:
\[
b'(s) \propto b(s)\prod_i P(e_i\mid s).
\]
The actions are ACCEPT, INVESTIGATE, and REJECT. INVESTIGATE obtains an additional independent probe and is capped at two rounds.

\section{Probability Model and Decision Rule}
The action model is cost-sensitive. False acceptance is assigned the highest provisional cost, followed by false rejection and then investigation. Direct actions compare expected false-accept and false-reject loss. Investigation is selected when its estimated one-step value of information exceeds its fixed cost.

\section{Test Method}
We generate 50 labeled simulation cases with five task categories. The cases include stable behavior, persistent degradation, transient anomalies, and distribution shift. We compare a static quality threshold, a binary policy, and the active policy. All cases use a fixed seed.

\section{Results}
See the generated table in \texttt{results/policy\_metrics.csv}. The experiment measures false accepts, false rejects, investigation rate, decision cost, degradation precision, and recall.

\section{Failure Analysis}
The active policy's highest-cost errors are exported to \texttt{results/five\_failure\_cases.csv}. Typical failure modes include weak evidence concentration, ambiguous transient signals, and distribution-shift cases that preserve average quality.

\section{Limitations, Ethics, and Human Control}
The likelihoods and costs are hypotheses, not calibrated production statistics. The synthetic cases are not historical incidents. A real deployment must validate judge reliability, calibrate thresholds, use comparable historical traces, and keep high-impact decisions under human control.

\section{Conclusion and New Questions}
The prototype tests whether active evidence collection can reduce false acceptance without making investigation excessive. The next experiment should replay public historical LLM generations and compare monitor decisions with human-reviewed labels.

\section*{AI-use statement}
AI tools were used to structure the research workflow, identify design gaps, and assist with code drafting/review. The project owner is responsible for the selected problem, evidence checks, experiments, interpretation, and public claims.

\bibliographystyle{plain}
\bibliography{references}
\end{document}
